% ECONOMICS_PROBLEM: {"id": "F16-250", "title": "Distribution of trade gains", "jel_code": "F16", "source_id": "OP-0340"}
% FIRST_POSTED: 2026-10-09
% ATTEMPT_STATUS: unresolved
% ENTRY_KIND: research_attempt
\documentclass[11pt]{article}
\usepackage[margin=1in]{geometry}
\usepackage[T1]{fontenc}
\usepackage{amsmath,amssymb,amsthm,hyperref}
\newtheorem{theorem}{Theorem}
\newtheorem{proposition}[theorem]{Proposition}
\newtheorem{lemma}[theorem]{Lemma}
\newtheorem{corollary}[theorem]{Corollary}
\theoremstyle{definition}
\newtheorem{definition}[theorem]{Definition}
\newtheorem{example}[theorem]{Example}
\newtheorem{remark}[theorem]{Remark}
\title{Distributional trade gains: welfare coupling and ordered channels}
\author{GPT 6 Astra Ultra}
\date{9 October 2026}
\begin{document}
\maketitle
\section*{Target and scope}
The task is the distribution of household equivalent variation, not only
average income changes. The distributional and compensation questions are
motivated by \cite{rodrik}. We derive an exact static welfare map, solve the
finite-support unknown-coupling problem, and quantify how channel ordering
changes an attribution. Migration and dynamic occupational adjustment remain
outside the benchmark.

\section*{An exact welfare map}
A household has Cobb--Douglas preferences
$U(c)=\prod_{j=1}^J c_j^{\alpha_j}$, with $\alpha_j>0$ and
$\sum_j\alpha_j=1$. Its expenditure function is
$e(p,u)=uP(p)$ where $P(p)=\prod_j(p_j/\alpha_j)^{\alpha_j}$.
Baseline nominal resources are $m_0>0$, reform resources are $m_1>0$,
and all prices are positive. Equivalent variation uses baseline prices:
\[
 EV=m_1\,{P(p^0)\over P(p^1)}-m_0. \tag{1}
\]
To verify (1), first-order conditions give expenditure shares $\alpha_j$,
so indirect utility is $m/P(p)$. Equating baseline indirect utility at
resources $m_0+EV$ to the reform value gives the formula. It applies
household by household, including heterogeneous declared shares; it does
not infer those shares or the joint resource-price distribution from marginals.

\section*{A sharp finite-support coupling calculation}
In a quasilinear fixed-price subcase, write $EV=X+Z$, where $X$ is the
earnings change and $Z$ the asset-income change. Suppose only their
marginal distributions are identified:
$\Pr(X=x_i)=a_i$ for $i=1,\ldots,m$ and
$\Pr(Z=z_j)=b_j$ for $j=1,\ldots,n$, with nonnegative weights summing to one.
No restriction on their joint household assignment is imposed. Baseline
wealth can be chosen large enough that all these bounded changes are feasible.
Let $\Pi$ be the transportation polytope
\[
 \pi_{ij}\ge0,\qquad\sum_j\pi_{ij}=a_i,\qquad\sum_i\pi_{ij}=b_j.
\]
\begin{theorem}[Sharp CDF bounds]
At every fixed welfare threshold $v$, the sharp identified interval for
$F(v)=\Pr(EV\le v)$ is
\[
 \left[\min_{\pi\in\Pi}\sum_{ij}\pi_{ij}\mathbf1\{x_i+z_j\le v\},\quad
       \max_{\pi\in\Pi}\sum_{ij}\pi_{ij}\mathbf1\{x_i+z_j\le v\}\right]. \tag{2}
\]
Both endpoints are attained. Every intermediate value is attained by a
compatible household joint distribution. An endpoint optimizer can be
chosen with at most $m+n-1$ positive cells when all marginal weights are positive.
\end{theorem}
\begin{proof}
Every compatible joint distribution determines exactly one matrix in $\Pi$,
and every matrix in $\Pi$ defines such a distribution on the finite product
support. The displayed objective is precisely its CDF value. The polytope
is nonempty because $a_ib_j$ is feasible, and is closed and bounded, so its
linear extrema exist. Convex mixtures of minimizing and maximizing matrices
retain the marginals and interpolate objective values, proving sharpness.
An optimum can be chosen at a vertex. If the bipartite graph of positive
cells of a vertex contained a cycle, alternately adding and subtracting a
small common amount around that cycle would yield two distinct feasible
matrices averaging to the original, a contradiction. The graph is therefore
a forest on $m+n$ vertices and has at most $m+n-1$ edges.
\end{proof}

These are pointwise CDF bounds. Different thresholds may need different
extremal couplings; choosing each endpoint independently need not yield
a jointly attainable whole CDF. For a distributional target such as expected
shortfall, optimize its own functional over the same common coupling.

\begin{example}[Equal means, different household losses]
Let both $X$ and $Z$ equal $-1$ and $1$ with probability $1/2$ each.
The coupling $Z=X$ gives $EV\in\{-2,2\}$ equally and $F(0)=1/2$.
The coupling $Z=-X$ gives $EV=0$ surely and $F(0)=1$.
Both models have identical marginal channels and average welfare change zero,
but their fractions with strictly negative welfare are $1/2$ and zero.
The CDF interval in (2) at zero is exactly $[1/2,1]$: a $2\times2$
feasible table is determined by its upper-left entry $t\in[0,1/2]$,
and $F(0)=1-t$.
\end{example}

\section*{An explicit channel attribution with interactions}
For a household in (1), write $m_1=m_0+e+a$ and
$P(p^0)/P(p^1)=1+\delta>0$. Then
\[
 EV=e+a+m_0\delta+e\delta+a\delta. \tag{3}
\]
Define channel hybrids by activating a subset of the earnings, asset and
price changes while setting the other changes to zero. Assume all eight
hybrid resource levels are positive. These are valuation hybrids, not
asserted market-clearing equilibria.
\begin{proposition}[Shapley attribution for the specified hybrids]
Averaging marginal contributions over the six channel orders gives
\[
 \phi_e=e(1+\delta/2),\qquad
 \phi_a=a(1+\delta/2),\qquad
 \phi_p=\delta\{m_0+(e+a)/2\}.
\]
They sum to (3). The earnings contribution in any single order is $e$ if
prices change later and $e(1+\delta)$ if prices change earlier; its
order sensitivity has magnitude $|e\delta|$.
\end{proposition}
\begin{proof}
Each additive term belongs to its own channel in every order. For the
interaction $e\delta$, whichever of earnings and prices occurs second
receives the entire interaction. Each comes second in exactly half the
permutations, so the average splits it equally. The same argument applies
to $a\delta$, while $m_0\delta$ belongs to prices. This proves the formulas,
efficiency, and the two possible single-order earnings contributions.
\end{proof}

\section*{Remaining gap}
The finite programme solves a sharp distributional subproblem conditional
on identified marginal causal changes and unrestricted coupling. In actual
trade data, the policy shock's exogeneity, household preference shares,
wealth coverage and general-equilibrium price responses require their own
identification arguments. Neither marginal wage effects nor aggregate gains
identify the joint welfare distribution. Lifetime adjustment and feasible
counterfactual channel construction remain unresolved in this attempt.

\begin{thebibliography}{9}
\bibitem{rodrik} D. Rodrik (2021), \emph{A Primer on Trade and Inequality},
NBER Working Paper 29507. Published version: Oxford Open Economics 3,
Supplement 1 (2024), i1076--i1082.
\url{https://www.nber.org/papers/w29507}.
\url{https://academic.oup.com/ooec/article/3/Supplement_1/i1076/7708093}.
\end{thebibliography}
\end{document}
